\chapter{Durability in Practice}
\label{ch:durability}

ARIES is the general answer for a page-based engine. A simpler engine can make a
simpler choice and still be correct. OntoDB's persistent databases, which run
today, are such a design: an in-memory database made durable by a
\textbf{snapshot} plus a \textbf{logical, redo-only write-ahead log}. This
chapter explains the design through the ideas of the previous chapters, because
every choice in it is one of those ideas in a smaller form.

\section{Where it sits in the steal/force table}

The database lives entirely in memory; the files are its durable copy. Nothing
is ever written to the data files except a whole snapshot, and a snapshot is
taken only between statements, when every change in memory is committed. In the
terms of \cref{tab:stealforce}:

\begin{itemize}
\item \textbf{No steal.} Uncommitted changes never reach the snapshot. So recovery never needs to \emph{undo} anything on disk.
\item \textbf{No force.} A commit does not rewrite the snapshot; it appends to the log. So recovery must \emph{redo} committed changes made since the last snapshot.
\end{itemize}

No steal plus no force is the \emph{redo-only} cell: the log needs redo
information and nothing else, and recovery is a single forward pass.

\begin{keyidea}
Choosing a cell of the steal/force table chooses the recovery algorithm. Steal
and no force need undo and redo (ARIES). No steal and no force need only redo.
The price of no steal is that all uncommitted state must fit in memory, which an
in-memory database accepts by definition.
\end{keyidea}

\section{Logical logging}

Page-based ARIES logs \emph{physical} or \emph{physiological} changes: this page,
these bytes, or this page, this slot. The memory database has no pages, so it
logs \emph{logical} operations: ``insert triple (18, 9, 4)''. Logical records are
small and independent of storage layout, and they are safe here only because
there is no partially written page that could make a logical operation
ambiguous on replay.

Ids in those records must mean the same terms after a restart, so the log also
carries the dictionary: a \term{term record} for each spelling the first time a
transaction uses it.

\section{One statement, one transaction}

\begin{enumerate}
\item The statement runs against a \term{recording store}, a decorator around the real store that remembers every triple it actually changed.
\item The log receives \textbf{begin}, a \textbf{term} record for each new spelling, one \textbf{update} per changed triple, and \textbf{commit}.
\item The commit is forced (\cref{ch:wal}) before the shell prints the result.
\end{enumerate}

If the statement fails part way, the recording store's changes are undone in
memory, in reverse order, and nothing it needs reaches the log. A statement that
changed nothing logs nothing.

\begin{figure}[ht]
\centering
\begin{tikzpicture}[font=\scriptsize,rec/.style={draw,minimum height=0.55cm,inner sep=3pt,fill=white}]
\node[rec,fill=boxbg] (b) {begin T2};
\node[rec,right=0pt of b] (t) {term 18 = \textless...\#Ceramic\textgreater};
\node[rec,right=0pt of t] (u) {update insert (18, 9, 4)};
\node[rec,right=0pt of u,fill=statusbg] (c) {commit T2};
\foreach \n/\l in {b/34,t/35,u/36,c/37}{\node[below=1pt of \n,font=\tiny] {LSN \l};}
\draw[arrow] ([yshift=6pt]c.north) -- node[right]{fsync, then ``inserted 1 triple''} ++(0,0.6);
\end{tikzpicture}
\caption{The log records of one \code{INSERT DATA} in a persistent database.}
\end{figure}

\section{Recovery}

Opening a database directory:

\begin{enumerate}
\item Load the snapshot, if there is one: every term in id order, then every triple.
\item Scan the log and stop at the first record that is torn or fails its checksum (\cref{ch:wal}).
\item Skip records at or below the snapshot's LSN; the snapshot already contains them.
\item Apply \emph{every} term record in order, whatever its transaction did: the dictionary only grows, and a term no committed triple uses is invisible.
\item Apply update records only for transactions that have a commit record, in log order.
\end{enumerate}

Step 5 is the whole of ``undo'' in a redo-only design: an uncommitted
transaction's records are simply never applied.

\section{Checkpoints as atomic snapshots}

A checkpoint writes the full state to \code{snapshot.tmp}, \code{fsync}s it,
\code{rename}s it over \code{snapshot}, \code{fsync}s the directory, and then
empties the log. Each step is safe to crash after:

\begin{center}
\begin{tabular}{@{}ll@{}}
\toprule
Crash after & On reopen \\
\midrule
writing \code{snapshot.tmp} & old snapshot, full log: nothing lost \\
the rename & new snapshot, old log: records at or below its LSN are skipped \\
emptying the log & new snapshot, empty log \\
\bottomrule
\end{tabular}
\end{center}

This is the log-structured idea of \cref{ch:lsm} with a single level: an
append-only log, periodically compacted into one sorted image.

\section{Damaged logs: recovery modes}

When the scan in step 2 stops early, there are two defensible policies, named
after RocksDB's WAL recovery modes:

\begin{description}[leftmargin=1.2em]
\item[Point in time] (the default). Keep every transaction before the bad record, truncate the log there, and report how many bytes were discarded. A torn final record is the normal shape of a crash, and this is the right answer for it; a transaction whose commit is after the damage is treated as uncommitted, exactly as a crash would.
\item[Absolute consistency] (\code{--strict}). Refuse to open, so an operator can inspect the files before anything is truncated.
\end{description}

A damaged snapshot always refuses to open: it has no valid prefix to fall back
to.

\section{Verification}

Detecting damage early is the other half of durability. The \code{\textbackslash
verify} command follows the tradition of PostgreSQL's \code{pg\_amcheck},
MySQL's \code{CHECK TABLE}, and RocksDB's \code{VerifyChecksum}: it changes
nothing and reports four checks, one line each.

\begin{lstlisting}
snapshot: ok (lsn 37, 18 terms, 15 triples)
wal: ok (0 records, 0 bytes)
replay: matches the live database (15 triples)
invariants: ok (no duplicate triples, every id in the dictionary)
verify: ok
\end{lstlisting}

The replay check is the strongest: it rebuilds the database from the two files
into a fresh store and compares it with the one in memory. If they differ, the
files would not reproduce what the user sees.

\section{Testing a recovery design}

A recovery path that is never exercised does not work. OntoDB tests this one with
\term{crash points}: named places in the code, such as just before a commit
record is appended and just after a snapshot is renamed, where a test can make
the operation stop as if the process died. Scenario tests then reopen the
directory and check what survived: a crash before commit leaves no trace, a crash
after the snapshot rename does not replay records twice, a torn tail keeps the
committed prefix, a flipped byte in the middle of the log drops everything after
it, and a tampered log fails verification.

\begin{ontodb}[: what runs]
All of this chapter runs today: start the shell as \code{ontodb mydb}.
\code{\textbackslash load} and every update are durable transactions,
\code{\textbackslash checkpoint} writes a snapshot, \code{\textbackslash log}
prints the decoded log, \code{\textbackslash verify} checks the files, and
\code{\textbackslash crash} exits without flushing, to try recovery by hand.
Seventeen scenario tests cover the cases above, and they pass under
AddressSanitizer, UndefinedBehaviorSanitizer, and ThreadSanitizer.

When the indexed store arrives, the same log records, log manager, and group
commit carry over; what changes is that pages can be stolen, so recovery gains
undo, pageLSNs, and the three passes of \cref{ch:aries}.
\end{ontodb}

\begin{decision}{Terms replay whatever their transaction did}
\dpart{Choice} On recovery, every term record is applied in order, even if its
transaction never committed; update records apply only for committed
transactions.
\dpart{Alternatives} Replay terms only for committed transactions.
\dpart{Why} Ids are dense, and each term record carries the id it was given.
If an aborted transaction's term were skipped, every later term would land on a
different id than the log recorded. A term no committed triple uses is
invisible.
\dpart{Cost} Aborted transactions can leave unused terms behind.
\end{decision}

\begin{decision}{Point-in-time recovery by default}
\dpart{Choice} A damaged log is truncated at the first bad record, and the
database opens with everything before it, reporting what was dropped.
\code{--strict} refuses to open instead.
\dpart{Alternatives} Strict by default; or skip the bad record and continue.
\dpart{Why} A torn final record is the ordinary result of a crash, and the
right response is to open. Skipping a bad record would apply later
transactions on top of a missing one.
\dpart{Cost} Damage in the middle of the log silently loses every later commit
apart from the report, which is why \code{\textbackslash verify} exists.
\end{decision}

\begin{checkyourself}
\item Why can the memory database's log omit undo information?
\item A crash happens after the snapshot rename but before the log is emptied. Which records does reopen apply?
\item A transaction logs a term record for id 19 and then fails. A later committed transaction logs id 20. Why must replay still apply the term for 19?
\item When would you choose \code{--strict} over point-in-time recovery?
\item Why is the replay check of \code{\textbackslash verify} stronger than checking checksums alone?
\end{checkyourself}

\section*{Further reading}
Gray and Reuter \cite{grayreuter} on redo-only logging; DeWitt et al.\
\cite{dewitt1984} on recovery for main-memory databases; the RocksDB WAL
recovery-mode documentation \cite{rocksdbwal}.
